\documentclass[UTF8,11pt,a4paper]{ctexart}
\usepackage[a4paper,margin=25mm]{geometry}
\usepackage{amsmath,amssymb,amsthm,booktabs,microtype}
\usepackage[hidelinks]{hyperref}

\newtheorem{theorem}{定理}
\newtheorem{lemma}{引理}

\title{行列式概率测度熵非凹性的一个六点对称反例}
\author{RenZhenzhuo\\\small 谨代表欧拉智能体团队}
\date{审阅稿，2026 年 9 月 6 日}

\begin{document}
\maketitle

\begin{abstract}
Lyons 曾问：有限行列式概率测度的 Shannon 熵是否是其正压缩核的凹函数。
本文给出一个定义在六个标号点上的反例。反例的中点是由秩三投影向内部
移动所得的实对称矩阵；两个端点则是在中点上加入一对符号相反的纯虚
Hermitian 扰动。借助对称性，64 个完备事件的概率可归为八类。一个简短的
严格计算表明，两个端点给出相同的概率分布，而它们的熵严格大于中点处的熵。
\end{abstract}

\section{反例}

令 $E=[6]$。对满足 $0\preceq L\preceq I$ 的 Hermitian 矩阵 $L$，
记 $X_L$ 为满足
\[
 \Pr(T\subseteq X_L)=\det L[T]\qquad(T\subseteq E)
\]
的行列式点过程，并记
\[
 p_A(L)=\Pr(X_L=A),\qquad
 H(L)=-\sum_{A\subseteq E}p_A(L)\log p_A(L).
\]
Lyons 在 2003 年提出了 $H$ 关于 $L$ 的凹性猜想~\cite{Lyons2003DPM}，
并在其 ICM 综述的猜想 2.6 中再次陈述~\cite{Lyons2014}。

\begin{theorem}\label{thm:main-zh}
存在严格正压缩算子 $Q_+$ 和 $Q_-$，其共同中点 $K$ 满足
\[
 p_A(Q_+)=p_A(Q_-)\qquad(A\subseteq E)
\]
以及
\begin{equation}\label{eq:gap-zh}
 \frac1{40\,000\,000}
 <H(Q_+)-H(K)
 <\frac1{39\,000\,000}.
\end{equation}
因而
\[
 H(K)<\frac{H(Q_-)+H(Q_+)}2.
\]
所以，行列式熵在复 Hermitian 正压缩算子集合上不是凹函数。
\end{theorem}

下面写出这些矩阵。令
\[
 R=\frac13
 \begin{pmatrix}
  2&-1& 2\\
 -1& 2& 2\\
  2& 2&-1
 \end{pmatrix},
 \qquad
 S=\begin{pmatrix}
 0&1&-1\\
 -1&0&1\\
 1&-1&0
 \end{pmatrix},
\]
并定义
\begin{equation}\label{eq:PB-zh}
 P=\frac12\begin{pmatrix}I_3&R\\ R&I_3\end{pmatrix},
 \qquad
 B=\begin{pmatrix}S&0\\0&-S\end{pmatrix}.
\end{equation}
最后取
\begin{equation}\label{eq:kernels-zh}
 \varepsilon=\frac1{475},\qquad
 t=\frac1{2300},\qquad
 K=\varepsilon I_6+(1-2\varepsilon)P,
 \qquad Q_\pm=K\pm itB.
\end{equation}
数值上，
\[
 H(K)=2.877726131217012\ldots,\qquad
 H(Q_+)=2.877726156388836\ldots.
\]
熵的增加量为
$2.517182345129681\ldots\times10^{-8}$，约为中点熵的
$8.75\times10^{-9}$。

\section{矩阵结构与端点分布}

\begin{lemma}\label{lem:structure-zh}
式 \eqref{eq:PB-zh} 中的矩阵满足
\[
 R^T=R,\qquad R^2=I_3,\qquad
 S^T=-S,\qquad RS=-SR.
\]
所以 $P$ 是秩三正交投影，$B^T=-B$，$PB=BP$，并且
$\lVert iB\rVert=\sqrt3$。
\end{lemma}

\begin{proof}
前四个等式由直接相乘得到。分块相乘随即给出
$P^T=P=P^2$、$B^T=-B$ 和 $PB=BP$。矩阵 $iS$ 的特征多项式为
$\lambda(\lambda^2-3)$，因此 $\lVert iB\rVert=\sqrt3$。
\end{proof}

$K$ 的特征值为 $\varepsilon$ 和 $1-\varepsilon$，重数均为三。由于
\[
 t\lVert iB\rVert=t\sqrt3<\varepsilon,
\]
特征值扰动界说明 $Q_\pm$ 的每个特征值都严格位于 $(0,1)$ 内。因此
$Q_\pm$ 都是严格正压缩算子。

下文使用完备事件公式
\begin{equation}\label{eq:event-zh}
 p_A(L)=(-1)^{|E\setminus A|}
 \det\bigl(L-I_{E\setminus A}\bigr),
\end{equation}
其中 $I_C$ 表示到坐标集 $C$ 上的对角投影。该式直接来自对包含事件概率
作容斥反演。

由于 $Q_-=\overline{Q_+}$，式 \eqref{eq:event-zh} 中相应的两个行列式
互为复共轭；而这两个矩阵都是 Hermitian 矩阵，故其行列式为实数。因此
\begin{equation}\label{eq:same-law-zh}
 p_A(Q_+)=p_A(Q_-)\qquad(A\subseteq E).
\end{equation}
端点分布相等这一点不需要任何数值计算。

\section{严格的熵比较}

令
\[
 u=\varepsilon(1-\varepsilon).
\]
保持 $P$ 不变并把 $B$ 送到 $B$ 或 $-B$ 的坐标置换，再加上取补集，
把 64 个子集分成下表中的八类。例如，坐标对称由
$(1\,2\,3)(4\,6\,5)$、$(2\,3)(4\,6)$ 和
$(1\,4)(2\,5)(3\,6)$ 生成。以
$\mathrm{diag}(I_3,-I_3)$ 作共轭会把 $Q_-$ 送到 $I-Q_+$；
结合 DPP 的补集规则和 \eqref{eq:same-law-zh}，一个子集与其补集具有
相同的中点和端点概率。将各类代表代入 \eqref{eq:event-zh}，得到表中
后两列。

先完整推导第一行。$K$ 的特征值 $\varepsilon$ 和 $1-\varepsilon$
各出现三次，所以
\[
 p_{\varnothing}(K)=\det(I-K)
 =\varepsilon^3(1-\varepsilon)^3=u^3.
\]
由 $PB=BP$，$P$ 的像空间和核空间在 $iB$ 下均保持不变。在
$x\mapsto(x,Rx)$ 与 $x\mapsto(x,-Rx)$ 的识别下，$iB$ 在这两个
三维子空间上的限制都由 $iS$ 表示，其特征值为
$0,\sqrt3,-\sqrt3$。因此
\begin{align*}
 p_{\varnothing}(Q_+)
 &=\det(I-Q_+)\\
 &=\varepsilon(1-\varepsilon)
   (\varepsilon^2-3t^2)((1-\varepsilon)^2-3t^2)\\
 &=u^3+3t^2u(3t^2+2u-1).
\end{align*}
这正是表中第一行。其余七行使用同一个行列式公式得到；随附检查器会逐一
核对所有代表及其重数。

\begin{center}
\setlength{\tabcolsep}{3.5pt}
\renewcommand{\arraystretch}{1.25}
\scriptsize
\begin{tabular}{@{}ccll@{}}
\toprule
代表 $A$&重数&$p_A(K)$&$p_A(Q_+)-p_A(K)$\\
\midrule
\(\varnothing\)&2
 &\(u^3\)
 &\(3t^2u(3t^2+2u-1)\)\\
\(\{1\}\)&12
 &\(u^2(1-2u)/2\)
 &\(-t^2(18t^2u-3t^2+12u^2-6u+1)/2\)\\
\(\{1,2\}\)&12
 &\(u(1-2u)^2/4\)
 &\(t^2(18t^2u-6t^2+12u^2-8u+1)/2\)\\
\(\{1,2,3\}\)&2
 &\((1-2u)^3/8\)
 &\(3t^2(1-2u)(3t^2+2u-1)/2\)\\
\(\{1,4\}\)&12
 &\(u(36u^2-20u+5)/36\)
 &\(t^2(54t^2u-12t^2+36u^2-10u+1)/6\)\\
\(\{2,4\}\)&6
 &\(u(9u^2-8u+2)/9\)
 &\(t^2(27t^2u-6t^2+18u^2-11u+2)/3\)\\
\(\{1,2,4\}\)&12
 &\((1-2u)(9u^2-4u+1)/18\)
 &\(-t^2(54t^2u-15t^2+36u^2-20u+2)/6\)\\
\(\{1,3,4\}\)&6
 &\((1-2u)(36u^2-4u+1)/72\)
 &\(-t^2(54t^2u-15t^2+36u^2-8u-1)/6\)\\
\bottomrule
\end{tabular}
\end{center}

下面说明如何在不依赖浮点数的情况下严格判定符号。将任意正有理数写成
$x=2^kr$，其中 $1\le r<2$，再令 $y=(r-1)/(r+1)$。对每个正整数
$m$ 都有
\begin{equation}\label{eq:log-bound-zh}
 2\sum_{j=0}^{m-1}\frac{y^{2j+1}}{2j+1}
 \le \log r\le
 2\sum_{j=0}^{m-1}\frac{y^{2j+1}}{2j+1}
 +\frac{2y^{2m+1}}{(2m+1)(1-y^2)}.
\end{equation}
在 $r=2$ 处应用同一个公式，可得到 $\log2$ 的有理上下界。因此
\eqref{eq:log-bound-zh} 为表中每个概率的对数给出定向有理界。

若某行的重数为 $m$，中点概率为 $p$，端点改变量为 $d$，则该行对熵差
的贡献是 $m[-(p+d)\log(p+d)+p\log p]$。代入
$\varepsilon=1/475$ 和 $t=1/2300$，将八行相加，并在
\eqref{eq:log-bound-zh} 中取 50 项，纯有理数运算给出
\[
 \frac1{40\,000\,000}
 <H(Q_+)-H(K)
 <\frac1{39\,000\,000}.
\]
这就证明了 \eqref{eq:gap-zh}，并由 \eqref{eq:same-law-zh} 完成
定理 \ref{thm:main-zh} 的证明。

随附的 \texttt{verify\_exact.py} 只使用 Python 标准库，从上述四个
矩阵出发，重建全部事件概率、八类分组、端点合法性、端点分布相等性以及
熵差的两个定向有理界。

\section{为什么这个方向有效}

上述计算还有一个直接的概念解释。若 $L$ 是实对称矩阵，$B$ 是实
反对称矩阵，则对 \eqref{eq:event-zh} 中的行列式作转置可得
\[
 p_A(L+isB)=p_A(L-isB).
\]
每个完备事件概率都是 $s$ 的偶函数，故它在 $s=0$ 处的一阶导数为零。
对熵求两次导数便得到
\begin{equation}\label{eq:second-derivative-zh}
 \left.\frac{d^2}{ds^2}H(L+isB)\right|_{s=0}
 =-\sum_{A\subseteq E}
 \left.\frac{d^2}{ds^2}p_A(L+isB)\right|_{s=0}\log p_A(L).
\end{equation}
通常为负、由一阶导数平方组成的项在这里消失了。因此，符号完全由事件
概率的二阶运动决定。用信息几何的语言说，这就是概率映射的第二基本形式
与熵梯度的配对。

这里的矩阵属于一个简单的代数族：只要 $R$ 是满足 $R^T=R$、
$R^2=I$ 的实矩阵，而 $S$ 是满足 $S^T=-S$、$RS=-SR$ 的实矩阵，
分块公式
\[
 P=\frac12\begin{pmatrix}I&R\\R&I\end{pmatrix},
 \qquad
 B=\begin{pmatrix}S&0\\0&-S\end{pmatrix}
\]
就自动给出正交投影 $P$ 以及与之可交换的反对称方向 $B$。上面的具体
有理矩阵取自这个族；选取它们，是因为相应事件概率只有八类，并且二阶
运动会提高熵。

这个族来自对正压缩算子集合边界附近一阶退化方向的搜索。该搜索解释了
构造的来源，却不是验证反例所必需的。给定 $R$ 和 $S$ 以后，证明只剩
引理 \ref{lem:structure-zh} 中的矩阵恒等式、完备事件公式
\eqref{eq:event-zh} 和八行熵计算。

本例的中点是实对称矩阵，方向 $iB$ 则是纯虚 Hermitian 矩阵。因此，
本例否定的是允许复 Hermitian 核的猜想；它不解决仅允许实对称核时的
独立问题，也不证明六是反例可能出现的最小维数。

本构造是在 Qwen Euler 智能体协助下找到的。数值搜索只用于定位结构；
一旦给定 $R$ 和 $S$，本文及检查器中的全部结论都由严格计算验证。
尽我们所知，该猜想此前没有公开发表的证明或反例。

\bibliographystyle{plain}
\bibliography{references}

\end{document}
